% ── sections/02_methodology.tex ──────────────────────────────────────────────

\section{Research Methodology}

\subsection{Datasets}

We use the six PPI network datasets from the EMOGI benchmark~\cite{emogi2021}:
CPDB (13,627 nodes), IRefIndex (17,013), IRefIndex-2015 (12,129), Multinet (14,398),
PCNet (19,781), and STRINGdb (13,179). Each network is stored as a dense adjacency matrix
in HDF5 format. To avoid materialising multi-gigabyte dense matrices (PCNet requires
$\approx$3~GB as float64), we implement a chunked sparse reading strategy: the adjacency
matrix is read in blocks of 2,000 rows, nonzero indices are extracted, and the result is
stored as a compressed \texttt{.npz} cache alongside the original file. Cached loading
reduces total data loading time from $\approx$96~s to $\approx$1.3~s for all six networks.

Node features are 64-dimensional multi-omics vectors comprising mutation frequency (MF),
DNA methylation (METH), gene expression (GE), and copy-number alteration (CNA), each
across 16 TCGA cancer types, with a fixed canonical ordering aligned across networks.
Cancer/non-cancer labels and train/val/test masks are loaded from the same HDF5 containers.

\subsection{EMGNN Baseline Architecture (Methodology 1)}

The EMGNN architecture~\cite{emgnn2023} operates in two stages, described in
Section~\ref{sec:theory}. A shared GNN encoder processes each PPI graph independently,
and the resulting node embeddings are aggregated through a meta-graph whose nodes represent
unique genes across all networks. A second GNN on the meta-graph produces final embeddings
that are classified by an MLP. We reproduce the benchmark using the original code with
GCN~\cite{gcn2017}, GAT~\cite{gat2018}, and GIN~\cite{gin2019} backbones, training with
Adam~\cite{kingma2014adam} (lr~=~$5\times10^{-3}$, weight decay~=~$5\times10^{-4}$),
2000~epochs, and early stopping with patience~250.

\subsection{Architectural Improvements (Methodology 2)}

We introduce four architectural modifications to the EMGNN baseline and evaluate them
via a controlled ablation study on CPDB with a GCN backbone.

\subsubsection{Residual Skip Connections}
Following He et al.~\cite{resnet2016}, we add element-wise residual connections between
consecutive GNN layers for all layers after the first. This improves gradient flow in deeper
configurations and stabilises training.

\subsubsection{Normalisation Strategies}
We optionally apply layer-wise normalisation after each GNN layer, controlled by the
\texttt{norm\_type} parameter with four options.

\textbf{BatchNorm1d}~\cite{batchnorm2015} normalises activations over the batch
dimension. Through ablation we find that, in the full-batch graph setting where the entire
graph is a single ``batch'', the running statistics are computed over all $N$ nodes
simultaneously, providing no meaningful distributional shift correction across training
steps. \texttt{BatchNorm1d} is therefore set as a configurable but \emph{cautioned}
option, disabled in our best-performing configuration.

\textbf{GraphNorm}~\cite{graphnorm2021} addresses this limitation by normalising
activations within each individual graph instance rather than across a batch. This preserves
meaningful statistics in the full-batch setting and is implemented as the recommended
alternative to BatchNorm for graph-level tasks.

\textbf{LayerNorm} normalises over the feature dimension for each node independently,
requiring no batch or graph structure information. This is the most general option and
operates correctly regardless of training batch size or graph partitioning.

Both GraphNorm and LayerNorm are implemented and exposed via \texttt{--norm\_type graph}
and \texttt{--norm\_type layer} respectively.

\subsubsection{Label Smoothing}
Hard binary targets are replaced with soft labels:
\begin{equation}
  \tilde{y} = y\,(1 - \varepsilon) + \varepsilon / K,
  \label{eq:labelsmoothing}
\end{equation}
where $\varepsilon = 0.05$ and $K = 2$. This reduces over-confidence on the imbalanced
cancer/non-cancer dataset and improves calibration.

\subsubsection{Learning-Rate Annealing and Gradient Clipping}
We use cosine annealing~\cite{cosine2017} to decay the learning rate from $\eta_0$ to
$\eta_{\min} = 10^{-5}$, and clip gradients to $\|\nabla\|_2 \le 1.0$ before each
optimiser step. These measures reduce late-training oscillation and prevent exploding
gradients in deep configurations.

\subsubsection{Feature Engineering Pipeline}
A preprocessing pipeline is applied to node features before training:
(i) variance-threshold feature selection (threshold~=~0.01) removes near-constant
features; (ii) Z-score standardisation centres each feature and normalises its variance.
The pipeline is fitted on the first network and applied consistently to all subsequent
networks to prevent data leakage.

\subsubsection{Hyperparameter Optimisation}
We use Bayesian optimisation via Optuna (TPE sampler, median pruner) to search over
learning rate, hidden dimension, number of layers, dropout, weight decay, normalisation
type, label smoothing, and LR scheduler. Fifty trials are run with 2000 training epochs
per trial and early stopping (patience~=~250). The best configuration is reported in
Section~\ref{sec:ablation}.

\subsection{Multi-Network Extension (Methodology 3)}

We extend the improved model to simultaneously train on multiple PPI networks and introduce
a learnable per-network importance weighting mechanism. A vector
$\boldsymbol{\theta} \in \mathbb{R}^{K}$ is learned jointly with the rest of the model;
at each forward pass, the weights $\mathbf{w} = \text{softmax}(\boldsymbol{\theta})$ scale
the embeddings of all nodes belonging to each network $k$, allowing the model to
up-weight high-quality networks automatically. We evaluate all six network combinations
reported in Table~\ref{tab:multinetwork}, ranging from two to six PPI databases, and train
for 2,000~epochs with patience~250 on an RTX~5090 GPU.

\subsection{Interpretability Analysis (Methodology 4)}

\subsubsection{Integrated Gradients Attribution}
We apply Integrated Gradients~\cite{ig2017} to attribute each prediction to individual
node features. Attributions are aggregated across all cancer-label meta-nodes by taking
the mean absolute value per feature, producing a global feature importance ranking across
the four omics modalities. The \texttt{AttributionAnalyzer} class provides a clean API
over Captum~\cite{captum2020}, with separate methods for node-level and gene-set-level
attribution.

\subsubsection{Gene Set Enrichment Analysis}
We apply two complementary GSEA methods over the MSigDB Hallmark 2020 gene set
library~\cite{msigdb2015}. First, Enrichr ORA~\cite{enrichr2013} is run on the top-200
predicted genes (by cancer probability), producing FDR-adjusted $p$-values. Second,
pre-ranked GSEA~\cite{gsea2005} uses the continuous prediction scores to rank all genes
and computes normalised enrichment scores (NES). Statistical significance is set at
FDR~$<$~0.05.

\subsection{Experimental Setup}

All experiments are run on an NVIDIA RTX~5090 GPU (CUDA~12.8, sm\_120) with PyTorch~2.7
and PyTorch Geometric~2.7. Random seed is fixed to 72 for reproducibility. The Adam
optimiser~\cite{kingma2014adam} is used throughout. Evaluation metrics are the Area Under
the Precision--Recall Curve (AUPR) and the Area Under the ROC Curve (AUROC), computed on
the held-out test set of the last specified network (CPDB in all multi-network experiments).
For single-network experiments we report both the best result across multiple independent
runs and the mean where applicable.
